\section{Awakening to Consciousness}

\begin{quotex}
In the beginning the world was nothing but the Atman, in the form of a man. It looked around and saw nothing different to itself. Then it cried out once, `It is I.' That is how the word `I' came to be. That is why even at the present day, if any one is called, he answers, `It is I,' and then recalls his other name, the one he bears.
\flright{\textsc{Brihadâranyata-Upanishad}}
\end{quotex}


In demonstrating the potential development of the increasing awareness of ``I'' consciousness, Evola builds on themes from the following books:

\begin{itemize}


\item \emph{Persuasion and Rhetoric}, by Carlo Michelstaedter


\item \emph{The World as Will and Representation}, by Arthur Schopenhauer


\item \emph{The Ego and its Own}, by Max Stirner


\item \emph{Sex and Character}, by Otto Weininger

\end{itemize}


The last named book was translated by Evola into Italian. There is one particular chapter --- \emph{The ``I'' Problem and Genius\footnote{\url{https://gornahoor.net/library/SexAndCharacter.pdf}}} --- that may be helpful before tackling Evola's short book on \emph{The Individual and the Becoming of the World}.

\paragraph{Stage One of Consciousness}


In Stage One, man is in what we could call the ``natural'' state, that is, he functions in a state of Nature. As such, he performs just as Nature asks: he breathes, eats, reproduces, and dies. As far as Nature is concerned, the Stage One man is complete. Whether or not man develops his Will or gains Wisdom, makes no difference to Nature. There is no gene to encourage higher states of consciousness, there is no natural process that will lead to the ``evolution'' of the spirit.

In Stage One, there is no firm sense of the ``I'', and the natural man is not autonomous. His instincts, his desires, his likes and dislikes, his opinions have their sources outside him and the I is passive in relation to them. For example, he won't eat Brussels sprouts because that is just ``the way I am''. Or if he desires something, he acts to fulfil it and does not question the desire itself. His opinions are absorbed from his milieu --- family, church, clubs, friends, political parties --- and serve more as a way to bond with his community or to maintain certain positions of power. His worldview comes from idle chitchat (Heidegger: ``\emph{gerede}``) with neighbours or commentary in the media.

Thus, for natural man, there is no burning issue of certainty. He seldom seeks to justify his opinions and beliefs. In debate, he is quick to begin with ``I think'' or ``I believe'' as though that were sufficient warrant for belief. He seldom sees a need to justify his opinion with facts or logic, nor to refute a contrary opinion. Even when through experience, he attempts such a justification, he argues as would a lawyer or sophist --- that is, simply to defend a point of view --- and not as a genuine seeker of knowledge. Nietzsche's dictum that ``the courage of one's convictions is not sufficient, what is necessary is the courage for an attack on one's convictions'' has no impact on him.

In this passage, Evola uses the word ``\emph{singolo}'' instead of ``\emph{individuo}``, both of which I have translated as individual (though I marked the former). A \emph{singolo} is like one stick in a pile, whereas an \emph{individuo} is an independent centre of awareness. The former type of individual is part of his environment and is not yet separated, that is, as a detached consciousness, the observer, Aristotle's unmoved mover. Such a one has no Will as such, what he calls ``free will'' is merely a spontaneity, or what Goethe called ``disorderly self will''.

\paragraph{Stage Two of Consciousness}


In stage two, man begins to awaken from the naive realism of stage one that posits a ``real'' world ``out there, right now'' based on nothing but his animal faith. Evola describes the attempts to account for this, first in the metaphysico-religious realm and then in science. Finally he ties this in to the quest for certainty and proposes how to ground this certainty.

In the next section, Evola proposes science as the solution to the awakening into Stage Two. The attempt is made to tame the chaos of phenomena through scientific laws. Yet, by the nature of the scientific method, even these laws are uncertain, since new theories always arise to replace earlier ones. Furthermore, science cannot answer the fundamental question: ``Why is there something rather than nothing?''

At this point, through self-awareness, the naive view of a ``subject'' knowing an ``object' is no longer tenable. Evola accepts Schopenhauer's view that phenomena are simply representations in consciousness and there is no ``real'' thing standing behind phenomena or causing them. Then the only certainty left is this self-awareness of the ``I''. This is the threshold to Stage Three.

% FALTA ALGO?

\subsection*{Privation}

Evola consolidates the understanding gained by the man at stage three with a discussion of the topic ``privation''. In Aristotle and Thomas Aquinas, ``privation'' is the absence of a given form in something capable of possessing it. As it is a lack, it has no being in itself, yet it is part of experience. Evola will use the concept of privation, which denies a real being, to refute the philosophy of realism, which postulates the being,

The realist claims that there is a something behind the appearances or representations that somehow causes or accounts for them. Evola denies this and reasserts his claim that the appearances are from the ``I''. But certainly there are things that are impervious to my will. Evola claims that they are not real things but rather a privation, something that reveals my own insufficiency. Along the way of explaining this, a new, and perhaps greater and more heroic, meaning is given to activity in the world.

In a note, Evola points out the limit at which philosophy as such must end; what is then required is self-realization, to which no study of philosophy can lead. Evola asserts that such a transformation is ``not a myth but a real possibility''. This cannot be overemphasized --- it is insufficient to study Evola, one must work to achieve such a realization by other means; otherwise, nothing will be understood.

The rejection of realism, the elevation of self-realization over rational thought, and the centrality of the will, freedom, and power have far reaching implications. This short work does not draw out all such implications, but it is an interesting exercise to ponder them. Evola writes:

\begin{quotex}
Therefore, with the intention to develop the position assumed by the consciousness of the third stage, we will consider the only true objection encountered by absolute idealism. In absolute idealism there is a doctrine that seeks to transform the negative task of criticism and scepticism that defines the second stage into something positive; and that ceasing to understand the world as a phenomenon, as simply an appearance, (the only legitimate conclusion of critical investigation) in order to understand it instead as something posited, created by the I. Therefore when one no longer speaks of representing but rather of positing and creating, the concept of a free will comes into play, and thus this problem arises: I can well reduce the world to my representation, but at what point can I reduce it also to my will and my freedom?
\end{quotex}


Since Evola's presentation is quite terse, it may be helpful to try to give some examples of the concept of ``Privation''. If we start with the assumption that the world is my representation and my will, then it is impossible to deny that the world does not seem to be subject to my will --- this is what Evola is referring to as ``privation''. Recall here what Guenon wrote about Man's role as Mediator between Heaven and Earth\footnote{\url{http://www.gornahoor.net/?p=1798}}. Heaven is the Idea. Man's role as Mediator is to manifest that idea in Nature. But nature is subject to its own law of Destiny, so it resists my will. That leaves us with two alternatives:

\begin{enumerate}


\item My inability to exercise my Will reveals in insufficiency in me... a \textbf{privation}. This is something I should possess, were I at the level of the True Man.

\item My inability to exercise my will reveals the existence of something to oppose it.
\end{enumerate}


The natural temptation, then, is option (2), to ascribe some independent reality as the true source or cause of my sense of lack. In other words, what is properly the will of the ``I'', gets projected onto something else.

\subparagraph{Projections and Privation}


We can look at three levels of projection:

\begin{enumerate}


\item The Persona projects itself onto the Shadow

\item The Mind projects onto the Body

\item The I projects onto the world.
\end{enumerate}

\subparagraph{The Shadow}


The man whose self is not well integrated experiences certain motivations, feelings, desires, attractions, weaknesses and so on as alien to him. Hence, he denies them in himself, and hides them in his unconscious. Jung called that the Shadow archetype. Notice the similarity here to what Evola is getting at: the Shadow is not a real entity, rather it represents a lack or privation in consciousness, something the I denies, yet it is a product of its own will. So ``I'' may project those feelings and desires onto others, or attribute them to outside forces like God, or the devil, etc. Oddly enough, someone else can easily see through ``my'' projections and recognizes them as my will.

\subparagraph{The Mind}


At the next level, the I sees itself as the mind alone (subtle manifestation) and separate from the body (gross manifestation). It experiences the needs and desires of the body as something alien that impinges on its consciousness. The body becomes that master of the ``I''.

\subparagraph{The Self}


Even if the ``I'' is integrated with the body, at a further level the body-mind organism separates from its environment. It then regards the world as an independent reality. The world seems dark and alien and is an obstacle to my purposes. Instead, I can understand the obstacle not as the ``other'', but rather as the revelation of my own insufficiency in that regard. This gives my activity --- that is, since my efforts to exert my will and bring such obstacles under control --- an heroic character. This is represented in many myths such as the labours of Hercules, Jason, the medieval search for the Grail, and so on.

\subsection*{True Will and Spontaneity}

If everything is the result of the I, the obvious question is, ``Why am I not always aware of creating the world?'' To address this question, Julius Evola makes the distinction between \textbf{spontaneity} and \textbf{Will}, which he first brought up in the discussion of the three

stages. By spontaneity, he means an immediate movement from the possible to the actual, while by Will, there is a ``space'' where many possibilities arise prior to one of them becoming actual. So when he says the world is ``my'' representation, it is in the sense of spontaneity --- that is, the world arises in my consciousness, though I do not choose it with full awareness.

That this is necessary follows from the argument thus far.

\begin{enumerate}


\item The Spirit has priority over the Soul, which has priority over the Body.

\item The Will is the unifying principle, bringing the activity of the Spirit into manifestation.
\end{enumerate}


Therefore, from point (1), The representation of the world in the Soul is created by the mind, whereas the common opinion is that the world creates the representations in the mind. So, from (2), this representation necessarily emanates from the Will. How does this relate to spontaneity?

As Guenon points out, Man is the mediator between Providence and the World. As Spirit, Man shares in Providence which is total Freedom. Yet as Body, Man is subject to the Destiny of the World Process. As Schopenhauer demonstrates, at that level Will is unfree, determinate, and blind. Tradition teaches that, over time, the World Process moves further and further from principial Unity; the material analog of this is the thesis that ``the entropy of the universe is increasing.''

Karl Marx wrote that the purpose of philosophy is to change the world. However, he did not mean change in the light of higher ideas, but rather in the sense of hurrying along the World Process in the direction away from Principle. Tradition, on the other hand, brings about change in conformity to the Will of God, or Providence''. Note that this is an actionless action, spontaneous but in a conscious sense. Julius Evola elaborates on this point:

\begin{quotex}
Here it is necessary to make a fundamental point and that is to understand the essential difference between spontaneity and will. There is spontaneity where the possible is identical to the real or to put it better, where that which is could only have been that, the act has the form of an inconvertible compulsion, of a brute happening and outburst, and it is passive, impotent in respect to itself. In the will, instead, there is an excess of the possible over the real, that is, one does not pass from the possible to the real immediately, but a point of autonomy, of ``\emph{potestas}``, controls the act as the final unconditioned reason for its being or not being, of its being this way or its being otherwise, as an act that is only one of the possibilities, or rather, compossibilities. It is important to note that spontaneity, as well as will, can be said to be free: however, while in spontaneity it is a question of a completely negative freedom, that is, of a freedom that simply means: ``not determined from something external''. In the will there is a positive freedom, that is, a freedoms that means absolute absence of conditions, whether they be internal or external, and therefore the contingency or, if one prefers, the arbitrariness of the act.

Once this distinction is understood, it is not so much based on concepts and intellectual subtleties, but rather on an immediate given of consciousness, on internal evidence that one either does or does not have; when the absolute idealist, in the face of the system of reality, asserts he was the ``I'' who posited it, it is evident that he is referring not to will, but to spontaneity. He is referring in fact to that activity whereby things are perceived and made interior to our I, to that fundamental ``assent'' so we become aware of them --- an assent that if it is a necessary condition for every reality, insofar as the reality experienced by the I (and we cannot coherently speak about other reality), it is still quite far from being a sufficient condition. In fact, in representing the real is not controlled by the possible, the I is passive in respect to its own act --- not as it asserts things but rather it is as if things were asserting themselves in it. As passion and emotion, the representation is something of mine, something that I pull from my own interiority (and at this point the legitimacy of the need for idealism arrives, the rest is satisfied until Leibniz), but it is not me, since I cannot give it freely to myself, since I don't stand in relationship of master to its determinations, so that the spectacle of reality is unfolding before me, that is, this reality, not the reality that I will. Consequently: to the extent that the idealist can say he was the I who ``posited'' nature, he reduces the I to nature, that is, in as much as that I, which is freedom, knows nothing, or, better said, acts as if it knows nothing, and, with an obvious paralogism, borrows the concept of the I from that of the principle of spontaneity. I can say it was I who posited nature, but in so far as I am spontaneity, not as I am properly an ``I'', that is, freedom and control.
\end{quotex}

\subsection*{Transformation is a Real Possibility}

Opposed to the Evola's view is philosophical realism, which Evola next tries to refute. Things have two aspects in regard to the individual, and these can be verified in one's own consciousness:

\begin{enumerate}


\item A thing is represented in consciousness

\item The I is impotent in the face of these representations
\end{enumerate}


Up to here, the magical idealist and the realist agree; they disagree about the nature of the impotency.

\begin{itemize}


\item Evola asserts that things represent an imperfect or self-limiting activity. Again, this is \emph{privation}, a quality the ``I'' should have by nature, but lacks, either because of an imperfection or a self-limitation. So, in conformity with Aristotle and Aquinas, this privation is non-being.

\item The realist cannot accept an imperfect or self-limiting activity; instead, he assumes a ``real'' cause, external to the I, to account for this impotency.
\end{itemize}


\begin{itemize}


\item Evola asserts that the universal or absolute is what comes first.

\item The realist believes that the particular things are what comes first.
\end{itemize}


\begin{itemize}


\item Evola claims that imperfection can be overcome by a strengthening of the Will.

\item The realist believes that that imperfection can be overcome by some process of evolution or development.
\end{itemize}


In refuting realism, Evola will hold that

\begin{itemize}


\item The I can be directly experienced as energy

\item There is no need to postulate something other than the I to account for limitations
\end{itemize}


Bear in mind that Evola is here concerned with the problem of certainty: that is, how do I know --- that is directly --- with certainty. For this, Evola only allows in evidence what happens in consciousness, my self-awareness or sense of ``I'', and, derivatively, what I have power over and can control.

Therefore, a ``thing'' is whatever appears in my consciousness as ``representation''. The realist, on the other hand, speculates that there is something that accounts for that appearance and thus causes it in my consciousness. Since what I can know with certainty is whatever is in my own consciousness, the theory of the realist must necessarily be speculative.

The objection of the realist is that you have no experience of directly causing or bringing about the appearances of the world. To counter this, Evola returns to the distinction between spontaneity and will that was made in the discussion of the stages. The appearance of the world is ``spontaneous'', in the sense that it arises of necessity without the full consent of my will. So whatever resists my will is not something that has an independent existence --- as the realist would claim --- but rather it represents a ``privation'', something that lacks being, and merely indicates an insufficiency on my part.

In Evola's own words:

\begin{quotex}
The realist, referring exactly to the point of real individuality, thus puts forward an appeal that is entirely legitimate. He puts us in front of a common occurrence of experience, a storm for example, and asks us if we can say that is was we who ``posited it''. While here the idealist would respond in the affirmative --- because as was said, for him ``to posit'' means simply to represent with ``free necessity''. We instead, referring to a positing that the principle of control and unconditioned freedom requires, would respond: ``That, actually, is not posited by the I.'' Then the realist immediately says: ``Since that is not posited by the I, there must be `something else' that posits it'' --- and infers a real or existent cause in itself of the representations, such as God, matter, the noumenon, etc. Instead, here lies the error and the point on which we are permitted to demand the full attention of the reader.

\emph{To say that I, as I, that is, as sufficient and free principle, cannot recognize myself as unconditioned cause of the representations, does not at all mean say that these representations are caused by ``something else'' and have some real or existent things in themselves as substrate, but means simply that I am insufficient for a part of my activity, which is still spontaneity, that such a part is not yet MORALIZED, that the I as freedom suffers a PRIVATION in it.}

Everything I cannot act on, everything that resists my will, is only a privation of this very will, something negative, not a being, but a non-being. Hence the realist must be rejected out of hand: he in his reference to ``something else'' --- God, noumenon, substance, etc. --- makes a being out of non-being, calls real something that, being only a privation of my power, being nothing other than a negation and a void in the unmultipliable body of my activity, should instead, by right, be called unreal. Thus he confirms this very privation --- as he eschews it: instead of the act that, by controlling and possessing them, annuls things and redeems the privation, he substitutes the act that recognizes them and superstitiously gives them being and an autonomous reality. To the former act, he instead applies the criterion of certainty of the third stage: that is, he demands that the free and naked I of the individual can genuinely assert the principle of absolute idealism, and however say: ``In truth, I myself am the cause and the Lord of this world, in which I live''. But when will it be possible to assert that? Obviously when the individual has redeemed the dark passion of the world in a body of freedom, when the form according to which he lives has let the representative activity (that is, that activity through which the spectacle of the universe is formed in him) pass from spontaneity --- from the harmony of the possible and real -- to bare, unconditioned causality, that is to: potent will.
\end{quotex}


Evola concluded his response to the realist with this claim:

\begin{quotex}
Obviously when the individual has redeemed the dark passion of the world in a body of freedom, when the form according to which he lives has let the representative activity (that is, that activity through which the spectacle of the universe is formed in him) pass from spontaneity --- from the harmony of the possible and real --- to bare, unconditioned causality, that is to: \textbf{powerful Will}.
\end{quotex}


At this point, philosophical arguments fall short and ultimately proof can only come from the development of one's will. Specifically, the appearances, representations, or ``things'' in the world do not arise spontaneously, but rather arise in conformity to one's will. In the text, Evola adds the following footnote describing the transformation from philosophical discussion to self-realization. I think it safe to conclude that this is not merely an abstract claim but rather represents Evola's actual experience.

\begin{quotex}
As this transformation, which we claim is not a myth but a real possibility, can then practically be completed, is a problem we treated --- at least within the limits in which it is possible to treat it publicly and generally --- elsewhere see \textbf{Doctrine of Awakening}, \textbf{Yoga of Power} , and so it is not found here. We can only say that it is a task which neither culture, nor devotion, nor philosophy, nor art, nor morality, nor anything other than what men call ``spirituality'', can make the least contribution. As to philosophy, its limit is magical idealism, in which it ends up by recognizing its own insufficiency and postulating the realization of power as that by which most problems can find their only absolute solution.
\end{quotex}


\subsection*{The Cosmic Value of Human Action}

Evola now asserts that the view of \textbf{magical idealism} reveals the \textbf{cosmic value} to human action. For the realist, everything already is, so that activity is little more than rearranging the furniture. Opposed to this view, action is heroic. The True Man creates being, where previously there had only been non-being or privation.

This perspective is difficult to grasp for anyone who hangs on to the tenets of realism, even if unconsciously. Evola is describing the True Will, the noumenal will, the unifying principle of the individual, transcendent to subtle and gross manifestation. The two views are antithetical.

\paragraph{Control}


\begin{itemize}


\item The realist can only understand this view as the self-affirmation of the ego with all that implies as far as control and domination over others for merely profane ends.

\item To the contrary, this is about control and domination over oneself to overcome one's own insufficiencies.
\end{itemize}


\paragraph{Lack}


\begin{itemize}


\item The realist experiences lack and tries to fill it from the things of the world.

\item Lack represents insufficiency and privation, to be overcome by creative action.
\end{itemize}


\paragraph{Eros}


\begin{itemize}


\item The realist takes as real the primal forces of \emph{eros} and   \emph{thumos} and seeks to satisfy them from the things in the world.


\item Eros and thumos are energies to be used to make my actions effective. As such, I control and dominate them to my own ends.
\end{itemize}


\paragraph{Fate}


\begin{itemize}


\item By his involvement with the things of the world, the realist submits his will to the Destiny of the World Process.

\item The True Will is Free and seeks to make itself manifest in the World Process.
\end{itemize}

This is how Evola expresses it:

\begin{quotex}
Now only in a view like this does the act of the individual have a \emph{cosmic value}, whereas in the view of realism every true meaning and purpose is removed from activity --- this should be clear to everyone. In fact activity truly has meaning and value only where it exists to make something real, that is not yet such. This case is verified precisely where the ``other'' --- or that which reflects the limit of my freedom --- is understood not as a reality, but rather as a negation and a void: then the world appears as something \emph{incomplete}, as something that demands its integration in that act of the individual, so that necessity is made freedom, to that development of self-assertion so that the potent actuality of the Unique one is laid out and re-asserted inasmuch as it is its privation.

If instead one posits that the ``other'' as such --- that is, exactly as that principle which limits my freedom --- is not a privation and a non-being, but rather a positivity and a reality --- then everything is already perfect, everything is already ``being'', and it is not necessary to do anything else. Every purpose and every value of activity and of becoming, every responsibility fails --- \emph{since the voids of my being are not also voids of being in general}: the ``other'', with the reality attributed to it, fills them up. Instead, in the other case the whole world appears as a dark, painful claim to the ``I'' so that one gives these to oneself according to power and, in that way, enacts it into being, redeems it from privation, and makes it real. And the becoming --- that which I make --- now has a value, a cosmic value.
\end{quotex}


\subsection*{The Creative Power of the Individual}

Evola concludes his discussion of privation with the observation that it is the Will that throws light on the unknown. The accumulation of knowledge is insufficient.

There is a continuum between spontaneity and free will. Everything objective and necessary in experience --- that is, everything experienced as ``privation'', arises from spontaneity. That means, there is no consciousness of having ``willed'' it; rather, it arises out of the Absolute as the experience of the potential becoming actual. From the perspective of the individual, it arrives ``out of the dark''. Evola elaborates:

\begin{quotex}
Hence the conception that appears in the third stage of the development of the individual is, as a whole, the following: a continuum of activity that has as limits spontaneity on one side and free will on the other. Spontaneity is the universal, free will is the individual. These limits are related to each other as potential to act: everything objective, immediate, and necessary in experience, is, in regards to the position of the individual, the non-being inherent in what is in potential --- and here perhaps we can understand what certain mystics alluded to when they spoke of the ``dark passion of the world'', the ``unspeakable suffering of existence'', in which the body of the ``celestial man'' is crucified.
\end{quotex}


With free will, this darkness is illumined. The ``act'' of the will makes the world ``real'', whereas privation is unreal.

\emph{The individual, now conscious of his power, is creative}.


This brings up three points

\begin{itemize}


\item The individual the source of the becoming of the world, thereby giving it its solidity and reality.

\item This is why the search for certainty in the things of the world will always end in frustration.

\item The individual, at this point, is the source of his own being.
\end{itemize}


He explains:

\begin{quotex}
Freedom is the act and the luminous flame of such a darkness, of such a privation; and the world becomes, it is made real in conformity with absolute reality only in and through this flame --- that is, only to the extent that the individual, asserting himself at the point of power and control, consumes, burns up his former nature that was spontaneously made.
\end{quotex}


Paradoxically, the so-called Traditionalist does not long for a return to the Past; rather it is only in creating the Future does he come to know his own powers and the Will of God Providence . Hence true explanations do not lie in things or in the past, but rather ahead, since they are dependent on the individual for their \emph{raison d'être}. Such an individual is his ``own property'' and ``persuaded/convinced''. (The latter two expressions are due to \textbf{Max Stirner} and \textbf{Carlo Michelstaedter} respectively.)


Evola then claims that it is the individual who gives meaning and purpose to the world. This distinguishes him from Schopenhauer's position that denies any such purpose; in the latter's system, the individual is not free and does not become aware of the noumenal ``Will'' as his own centre. (I think it should be clear by now that Evola is not exalting the ``human being'' as such, since the states he describes are beyond the human state.)

\begin{quotex}
only at the point at which the individual realises himself in the sudden intuition of power does a purpose, a reason and a goal in nature arise: not before; it is he who gives it to them. It demands it of his action. And yet the individual has a single imperative by this point.
\end{quotex}


Evola has been describing the Traditional esoteric path. It consists of three stages.

\begin{description}

\item[Catharsis] This is the stage of purification, moving from spontaneity to Will.
\item[Theoria] This is the stage of illumination, where the darkness is illumined.
\item[Theosis] This is the stage of divinization, where one incarnates God.
\end{description}


To some it may seem that Evola is inverting the process of salvation in the Christian sense. Others may recognize in it the Traditional teachings of the Greek fathers. Evola concludes:

\begin{center}\textit{Be, make yourself GOD, and in bringing that about, SAVE the world.}\end{center}

\section{Builder of Worlds}

\begin{quotex}
It is important to note the relativity of the concept of privation. A given element is never privation in itself, but always in relation to the good of autonomy. The passage to that good makes what was positive as spontaneity something negative and ``in potential'' in respect to the other point. So for the man who does not want to move from the logical point of view to that of Will, the concept of privation is not intelligible.
\end{quotex}


Let us try to illustrate Evola's point with a homely example. Suppose Charlie has a problem with drugs; he is powerless with respect to cocaine. He is not autonomous because the desire for the drug controls his actions. But that is relative to Charlie because Cologero has no such problem. Thus, it is not the cocaine as a positive force, but rather a privation, a lack within Charlie, that accounts for the addiction. So it is unconscious, ``spontaneous'' will that makes of cocaine a positive reality, seemingly outside Charlie. However, as Charlie develops a conscious Will and autonomy, that ``positive'' addiction is recognized as a negative, as a privation, something that can be overcome.

\begin{quotex}
By believing that the present doctrine explaining privation is surpassed by postulating a distinct reality, he is not taking a step forward but a step backward, since he is making use of the logical category of causality with which this very reality becomes conditioned, logically posited by the I.
\end{quotex}


The realist cannot resist the urge to postulate an external reality: Charlie's problem is genetic or something of that nature. This doesn't address the problem, but pushes it back into more obscurity. The realist forgets that causality is not a property of the World, but is a category of thought by which I organize and control the World.

\begin{quotex}
What is the difference between a real and an imagined thing? Represented, they are both the same; but beyond that, the representing activity to which the real thing corresponds is an activity in respect to which they are impotent. There are elements on which I cannot act... We do not resolve the problem of interpreting this non-power, because we do not pose it; hence we are accused of being intellectualistic, abstract, and irrelevant in respect to what is really important to every research of such a type at this point. This is a fundamental point: we claim that the explanation of the fact that one is impotent in certain situations by making recourse to an ``other'' is a pseudo-explanation, and therefore a vicious circle for this reason: in us the concept of ``other'' gets its meaning and its foundation from the concept of ``non-power'', which is what comes first, and out of which objectivity, thing in itself, God, etc. are only so many symbols and intellectual interpretations... \emph{Privation} explains the concept of objective reality but not objective reality the concept of privation.
\end{quotex}


So I have a representation of things in my mind and then note my lack of power of some things. I start with the awareness of my lack, and then postulate the existence of the ``other'', all claiming to ``explain'' my awareness of privation. The awareness is certain, but there are many conflicting theories about the ``other''. That is why Evola considers them to be pseudo-explanations.

\begin{quotex}
The explanation that \textbf{magical idealism} demands is completely different: it is an \emph{explanation by means of action, a resolutive explanation}. It is to \emph{ex-plicate}, or to \emph{actuate}, to make perfect: to make what is in potential pass into act, what is imperfection into perfection, what is insufficiency into sufficiency, according to a synthetic, creative, primordial process. This is the only true explanation. \emph{Everything else is a pastime}.
\end{quotex}


If knowing is the Unmoved Mover, according to Guenon, and given the fundamental metaphysical principle ``to know is to be'', then to ``know myself'' is to ``be myself''. To be myself, I must actualize myself. To be the True Man, I must actualize all my possibilities. There is no distinction here between knowing and doing, contemplation and action: knowing cannot exist apart from acting. It is only by doing can I know myself, through the reflection of the I in the World. By creating myself, I reveal myself. Quietism, or the surrender to impotency, is just a form of escapism.

\begin{quotex}
We will fiercely combat all intellectual and philosophic rhetoric whereby man restrains himself to talk about his impotency (which we mean when speaking of ``truth'', ``objectivity'', ``rationality'', etc.) rather than to finally jump to his feet, to grasp himself and, burning up his impotency, to make himself what he is in himself:

\begin{center}\textit{a God, a builder of the world.}\end{center}
\end{quotex}


\subsection*{Essence and Existence}

\begin{quotex}
The universe exists for each Individual because he consciously wills it.

\flright{\textsc{John Woodroffe}, \emph{Sakti and Sakta}}
\end{quotex}


After concluding the discussion of privation, Julius Evola next turns to the concepts of Essence and Existence. This can be understood as the philosophical analog of the concepts of \emph{Purusha} and
\emph{Prakriti}. There are some important notions developed in what Evola wrote about essence and existence. Evola avoids a one-sided philosophy. He relates these concepts to freedom and necessity, will and thought, the actual and the potential. Yet he accomplishes this not by an unbridgeable dichotomy between opposed concepts, but rather by relating them on a scale of values.

Evola offers us a full philosophy: there is transcendence, yet there is activity in the world. The first point he makes is that things comprise both essence and existence. The essence of something answers the question ``What it is''. Existence is ``that it is''. Yet the idea of a thing is logically the same both for an imagined thing and a real thing. Therefore, the existence of a thing is irreducible to logic. In practical terms, what this means is that there is no intellectual answer to the question: ``Why is there something rather than nothing?'' The solution to that riddle lies on a different level.

\begin{quotex}
Things are essence and existence. The idea of one hundred thalers and one hundred real thalers are obviously not the same thing. Consequently, from a logical point of view, as Kant has shown, there is nothing included in the hundred real thalers beyond the idea of the hundred thalers. It follows that, insofar as we differentiate between them, we are referring to something irreducible to the logical element. This something is ``existence'', as opposed to ``essence'' (or more rigorously, the ``esse existentiae'' as opposed to the ``\emph{esse essentiae}``).
\end{quotex}


The essence of a thing is logically explained by a concept. However, its existence is beyond logic (note: this does not mean illogical or irrational). Since in Evola's system there is only representation and the ``I'', therefore existence can only be explained by will or by power. Otherwise, the existence of a thing is opaque to the I. Of course, tying this back to previous topics, the activity of my will may be spontaneous, i.e., beyond my conscious intervention. Nevertheless, it ``explains'' existence without resorting to theories of external forces. So, here is the symmetry:

\begin{enumerate}


\item essence is the ideal and is constructed in thought

\item existence is the real and is caused by an act of will.
\end{enumerate}

\begin{quotex}
And now to a second point. The explicative principle of the essence, of the ``what is'', of a determinate reality, is the concept: when a reality is genetically constructed by means of the concept in all the characteristics that determine it, the explicative principle in the order of essence is exhausted. Therefore, what an object is, which is entirely penetrated by being, is the bare fact of its ``being there'' as real object, and that constitutes a point that entirely eludes a rational explanation, it is an \emph{alogos} --- and an adequate explicative principle for it is not the concept, but rather the \textbf{Will} or, even better, \emph{power}. In fact the pure being of things constitutes for me a mystery up to the point where it has the character of a brute given, of something that is there without the participation of my will, imposing itself, on the contrary, forcefully against it; briefly: as a privation of my activity. While I can think of essence and therefore ``construct it'', existence simply suffers it --- and this constitutes for me a darkness. If you imagine instead a situation in which you can connect the ``being here'' of things to willing them unconditionally, that is, in which my will had value as creative power: then their existence in fact beyond their concept would cease to be a mystery to me, on the contrary, it would be for me perfectly intelligible --- it would be explained. \emph{Essence and existence have therefore as their explicative principles the ideal construction through the act of thought and real causation through the act of the will respectively}.
\end{quotex}


Evola adds a third point. He rejects dualism, in particular a dualism between thought and action. The world is representation and will, essence and existence, thought and action. The third point is: there is not a difference of nature, but only of degree. Note how this is another expression of the Great Chain of Being\footnote{\url{http://en.wikipedia.org/wiki/Great_Chain_of_Being}}, that reality is hierarchically ordered from the less dense to the more dense. Thus the idea is already at a level in the chain and the real is a fuller expression of the idea. The Tantric writings, as documented by John Woodroffe, also portray a similar conception; they add many levels of detail in describing the full constitution of the human being, for example.

\begin{quotex}
The third point is the following: between ideal constructions and the creative will --- therefore between essence and existence --- \emph{there is not a difference of nature, but only of degree}. The idea is already a degree of real affirmation; and so-called objective reality is only the most intense and complete assertion of that power which, in elementary form, determines the thing simply thought or represented. Reality is only the \emph{act} of the idea, in which it characterises and expresses itself entirely, just as the idea is only a reality \emph{in potential} or a reality simply sketched out or in a nascent state. Between one and the other, therefore, there is not a jump, there is instead \emph{progressivity}. The thought of one hundred Thalers and one hundred real Thalers are obviously not the same thing --- but not \emph{qualitatively} different (as is assumed by those who believe that the thought is the impersonal image of an objective reality instead of a power) but intensively different, in the sense that the hundred real Thalers are the deepest, most \emph{intense} power related precisely to the magical act, of the assertion corresponding to the hundred Thalers in thought.
\end{quotex}


This has important implications. First of all, there is the first type of existence that is ``death, privation, unreality'' which is the spontaneity of the first stage. Life is experienced as impotence. Things and events block me. I impute reality to the unreal. Thinking deals only with concepts and I can never grasp the real simply by thinking about it. To escape, I must be convinced that philosophy is incomplete without a method of self-realisation.

Only in that way --- and not by ``reasoning'' oneself to it, can there be the second type of existence. This type is based on the freedom, power --- beyond logical analysis as was shown in the second point. Notice the comments about ever increasing levels of spontaneity. Thus as a man's power is expanded, his actions become effortless, i.e., they arise spontaneously from the nature of his being, the \emph{wei wu wei} of Taoism.

\begin{quotex}
There is an existence that is death, privation, unreality, which corresponds to representational spontaneity, the residue from the first stage, in which the act is passive in respect to itself, and which the ``I'' does not control as its master. There is no true certainty of this existence: since it does not depend on me as passion or emotion and is a pure happening, therefore a principle of radical contingency takes it back. There is instead a second existence, which is what a will elevating itself to power can unconditionally produce: only this is properly existence, absolute reality, and only from it, where it is found joined only with itself in possession and control --- the ``I'' can have real certainty. Between these types of existence there is mental activity in the proper sense. In other words: beyond the ideal limit of the reign of pure necessity --- of nature and spontaneity --- and beyond his ``privation'', the individual enjoys in the rational or ideal order of the first level of sufficient actuality and freedom. This level proceeds towards its perfection in development in accordance to which power reasserts itself in continually deeper and more complex levels of spontaneity --- of primitive nature or the universal --- right up to controlling the same intensive level of real existence. Then from the dark passion and savage desert resulting from privation, the world will make itself the very act of the individual, and in that he will be redeemed and \emph{convinced}.
\end{quotex}

\subsection*{Knowledge and Power}

\begin{quotex}
The fundamental being and expression [of Consciousness or the I] is Joy pulsating as Will-Power and manifesting Itself in an unspeakably sublime cosmic play. It is not a mere abstraction — a wilderness of Pure Being or Pure Nothing as some critics of Vedanta have imagined the abode of Reality to be. \flright{\textsc{John Woodroffe}, \emph{The World as Power} (p 320)}

\end{quotex}
Thomas Aquinas describes the effects of the Fall from the Primordial State as affecting faculties on both the spiritual and the psychic level, as represented in Table \ref{tab:FallPrimordialState}:

\begin{table}[h]
\centering\small
\begin{tabularx}{\textwidth}{llp{.3\textwidth}l}
\toprule
 &
\textbf{Effect} &
\textbf{Description} &
\textbf{Antidote}\\\midrule
Spirit &
Ignorance &
Intellect not ordered to the Logos &
Gnosis\\\cmidrule{2-4}
 &
Malice &
Will impotent to manifest the Logos &
Power\\\midrule
Soul &
Weakness &
Thumos not ordered to arduousness &
Trial by suffering\\\cmidrule{2-4}
 &
Concupiscence &
Eros not ordered to the delectable &
Trial by love\\\bottomrule
\end{tabularx}
\caption{Fall from the Primordial State}
\label{tab:FallPrimordialState}
\end{table}
Hence, the regeneration of the Primordial State is effected through Knowledge and Power.

\paragraph{Knowledge}
What do we know when we know it? From what has been described thus far:

\begin{itemize}
\item All our opinions, beliefs, theories are like so much straw 
\item What we have considered as obstacles are in reality privations, a demonstration of our own lack of power. 
\end{itemize}
Woodroffe describes the concept of privation as understood in Tantra:

\begin{quotex}
Consciousness [the I] \emph{seeming} to be unconsciousness, Joy \emph{seeming} to be indifference or pain, free Play \emph{seeming} to be necessity and determination. 

\end{quotex}
In this seemingness, the I gives up his power to that which ``seems to be" and doesn't recognize his own role in producing what he experiences:

\begin{quotex}
Consciousness, as the root Being-and-Becoming Power becomes real things … Thus, a block of stone is not ``matter" only: it is Consciousness as Joy and as Play — though the fact is veiled to ordinary minds. On the other hand, it is not an ``idea" or ``mental construct" only: it is Consciousness as Power constituting it as much and as active a reality as the experiencing and reacting mind is. 

\end{quotex}
The solution is, according to Woodroffe, increasing knowledge at ever deeper levels to the point of self-realization:

\begin{quotex}
There is need for the education and development of man's ``knowing instruments", giving him progressively higher and larger visions, through science and philosophy, through intuition and meditation, and finally revelation and realization. … \emph{Matter and Life are every whit as real and active as Mind}.

\end{quotex}
This is the crux of the issue. Verbosity, opinions, repetition do not and never will constitute knowledge. \textbf{To Know} — in the highest sense, viz., episteme, gnosis, wisdom — is \textbf{to be}. Hence, we do not know the Idea until it has become Real for us as Life and as Matter. This is the Tantric teaching of the I as Power.

\paragraph{Power}
In the \emph{Individual and the Becoming of the World}, Evola describes the deepening process of making ideas real:

\begin{enumerate}
\item First, there is the simple and pallid image of the idea 
\item Make it vivid in the imagination 
\item Perceive it exteriorly as a subjective hallucination 
\item Convince other centres of consciousness to perceive it as a collective hallucination 
\item Develop an ever more intense assertion of power until the same level of real existence is reached. 
\end{enumerate}
Then at this point:

\begin{quotex}
in correspondence with a magical act (because a \emph{magician could be defined as one who knows how to influence nature itself}), [the idea] would cease to be ``false" and ``subjective" by making itself instead ``true" and ``objective". What was error or illusion becomes accordingly truth, through power. That is valid for any empirical, moral, or rational field whatsoever. 

\end{quotex}
Evola here does not spell out the obvious corollary. Whatever happens in any ``empirical, moral or rational field whatsoever" is the result of the intention of an agent, someone or group that understands, and uses, this process or, better said, ``magicians". How much of the world is under the spell of some collective hallucination or another? Only the Trial by Fire can break the spell. Evola's goes on to demonstrate the relationship between ontology and ethics.

\begin{quotex}
Real means perfectly willed. Unreal will mean therefore privation of the Will, imperfect and insufficient Will — and this is \emph{evil}, just as, on the other hand, it is error. \emph{Evil and error are one and the same thing in the sense that the common foundation to both is impotence}. Those doctrines are in truth of a fundamental value, currently almost forgotten, that placed evil in matter, meaning by this that residue of necessity that resists the idea — that is, freedom — and limits it. Matter — the Platonic ``other" — is effectively the sign of imperfection of the act and, in as much existing, is connected to that fundamental ``injustice" about which Anaximander, Parmenides, and Empedocles speak — and it is, essentially, evil. There is no other evil beyond necessity, of which matter, brute existence, is the evidence; and as long as the ``I" will not be able avoid the experience of matter — of this ``other" against and beyond him — he will be imperfect: \emph{evil, impure, irrational}. 

\end{quotex}
It is clear that freedom is not about satisfying desires, which is really a lack, a privation. Concupiscence represents a necessity that binds and restricts our will. The purpose of the Trial by Love, by substituting the satisfaction of personal desires with a disinterested love for the genuinely good or pleasurable, is to overcome concupiscence. Thus, commands such as ``love your enemy" are to be understood in this light, as means to an end, not a universal moral principle in the Kantian sense. A blind hatred, on the other hand, is the result of an impure and irrational will, and often leads one to make poor or ineffective decisions.

Weakness of the \emph{thumos} leads a man to eschew suffering. Every form of pain or upset is felt to be an injustice. Instead, the Trial by Suffering will develop a detached attitude. In this respect, we can mention the idea of ``intentional suffering", that is, deliberately choosing the more arduous path in order to develop Power. Once again, commands such as ``turn the other cheek" must be understood in this light.

Now Evola summarizes his position on the identity of power and morals:

\begin{quotex}
So the ethical problem and the metaphysical problem coincide: the measure of reality, as of the good, of the certain and of the true, is the perfection of actuality or power. \emph{Per virtutem et potentiam idem intelligo}. [Spinoza, \emph{Ethics}, ``by virtue and power I mean the same thing."] The good or virtue is potency, evil is impotence. There is no other evil except insufficiency and weakness, no other good except the will that is unconditionally sufficient in itself, and therefore free to be and do what it wills. The rest is nonsense and can be left to women, mystics, and poets. That today the meaning of the moral as cosmic worth has been lost, that up to now, it has been reduced to nothing more than the corroboration and canonisation of deficiency, weakness, and fear — of the plague of beautiful sentiments, of noble virtues, of holy ideals — that is, of the fundamentally \emph{immoral}. 

\end{quotex}

\subsection*{The Religious Attitude}

\begin{quotex}
In a nutshell, this is the sense of the system I uphold, in which on the one hand I sought to fuse the gnoseological problem and the ontological problem with the ethical problem and the problem of self-realisation or magical problem; on the other hand, to claim the value of the individual and to make him become conscious of his task and his cosmic dignity.

\textbf{It is that which I recognize as truth, or, better said, it is that which I will as truth.}
\end{quotex}


Before addressing the topic of initiation, Evola first summarizes his philosophy. First, he describes the initial state of the being.

\begin{quotex}
we have in fact the simple state of the being that discovers itself, that is pure spontaneity, that does not \emph{own} itself but simply \emph{is}. Although a state of fullness and light for the ``I'' yet to be born, close to the point of the individual, it appears instead as darkness and death.
\end{quotex}


The being is not master of himself as an I, but passive. Initiation looks forbidding, but will be enlightenment for the twice-born. He recapitulates the three stages.

\begin{enumerate}


\item In a first moment, it is dissolved in the world of appearances and mere representations.

\item In a second moment it is felt as infinite passion, as the dark and silent pain of privation, as the indescribable crucifixion in the world of necessity.

\item But, born from it, the individual now assumes this death with joy; he is adequate to it; he knows that only his own supernatural value of `being made of possession'' is its cause; he recognizes it as matter, from which only he will be able to drag the splendour of a life to an absolute reality.
\end{enumerate}


Evola describes this process of initiation in more detail. The ``being of possession'' echoes back to Max Stirner; the being now ``owns'' himself. The ``three worlds'' is a reference to Taoism.

\begin{quotex}
And then the darkness gradually is illumined, then the \emph{terrible} flower of the absolute Individual rises from the abyss of necessity. He raises himself slowly into the starless sky, pulling himself from the fierce heat of what he devours in his power. Things and beings \emph{die} in his vertiginous intensity and he gradually, irresistibly, becomes --- he, \emph{tremendous} in his purity, is ``Master of Yes and No'' and Ruler of the ``three worlds''. And in him, a being of possession, a being that ``burns and blazes'', the process of the universe will have its consummation or final perfection with his act.
\end{quotex}


In the second part of \emph{The Individual and the Becoming of the World}, Evola distinguishes between the ``religious attitude'' and initiation. The former has these characteristics:

\begin{enumerate}


\item It believes that the world is fundamentally right from a principle of order, harmony, and goodness, and that everything opposed to it that is darkness, irrationality, indeterminateness, evil, either is illusion or something that must not be.

\item Consequently, it pushes such a principle back onto something ``other than me'' -- to something transcendent. We say ``consequently'' since my actual experience is showing me a world totally different from one that is ordered and rational, I am not entitled to also admit the actual existence, alongside it, of a providential and rational world, unless I relate it to something that transcends me and to believe that this something exists.

\item If such a transcendent principle is intended to be the source of every reality, it must be said that the individual in himself is nothing and can be nothing. The ``I'' is annulled and its being is pushed back onto the other.
\end{enumerate}


Evola contrasts that to the wisdom of the mysteries --- of Buddha, Dionysus, Mithras, Hermes:

\begin{enumerate}


\item He faces up to reality, without veils, in its tragic, irrational, unprovidential nature.

\item He rejects the belief in the present existence of an optimistic order, and the necessity of postulating the transcendental, the God to guarantee it, and to explain how it can coexist with the order of existing things that contradict it.

\item Thus the gulf remains between titanic pessimism and Dionysianism, bridged by recognizing or not recognizing the moral right to existence of a regulated, just, and rational world. In the case of the affirmative, he asserts the concept of the individual as a sufficient power in himself, capable of saving not only himself but also everything that he is connected to and that, at bottom, is the reflexion of what he concretely is.
\end{enumerate}


\flrightit{Posted on 2011-03-17 by Cologero }

\begin{center}* * *\end{center}

\begin{footnotesize}
\begin{sffamily}
\texttt{Zmaj on 2011-03-05 at 05:20 said:}

Yes, this is a true summary.

I mention two additional points:

1) Readers who are familiar with the subjects and authors discussed are likely aware of the distinction between `will' in the lower sense and the `Will' treated in this post.

The `will' in most people is dominant, tyrannical even, and it subjugates the intellect, reducing it to discursive thought. This lower will usurps a higher rank that does not belong to it. Those who are thus arranged (i.e. the average majority) will find the current decadent society quite accommodating; their will is constantly pulling them into manifold activity, and they only need to think about how to do those things that will satisfy their whims and schemes. Their internal rebellion parallels the external disorder.

The Will is of a different quality to `will'/desire. It is rightfully situated above or alongside the intellect, and acts to integrate and strengthen the intellect. When Will is awake, proper order is established, with the intellect governing the thoughts, whilst the clarified thoughts fortify the ordinary will; and this will not necessarily lead to a life of `high achievement' in the social sense. Simultaneously, the Will is outside of the realm of the persona, and this is its proper place -- the plane of real activity.

Transformation is indeed a real possibility, but a dangerous one. A severe trial can accompany an attempt, and success is not guaranteed. Veil thyself with the darkness of a valley at midnight before all else.

\hfill


\texttt{HOO on 2012-03-23 at 00:36 said:}

``The suffering from frustrated ambition is peculiar to people living in a society of equality under the law. It is not caused by equality under the law, but by the fact that in a society of equal\-ity under the law the inequality of men with regard to intellectual abilities, will power and application becomes visible. The gulf between what a man is and achieves and what he thinks of his own abilities and achievements is pitilessly revealed. Day\-dreams of a ``fair'' world which would treat him according to his ``real worth'' are the refuge of all those plagued by a lack of self-knowledge.''\\ --- ``The Anti-Capitalistic Mentality'' by Ludwig von Mises

\hfill

\texttt{Aidan on 2011-03-08 at 14:07 said:}

Supra-individual states are not associated with individuality. This is not a linguistic turn but a point that is mathematical in its precision. Entry into supra-individual states requires the complete surrender of all individuality -- the partial preservation and integration of specific elements of the latter within the former is a matter which is not in the hands of the individual.

``It is almost superfluous to stress how little place the individual `self' occupies in the totality of the being, since even given its entire extension when envisaged in its integrality, and not merely in one particular modality such as the corporeal, it constitutes only one state like the others, among an indefinitude of others.''

and

``... there is assuredly a form of consciousness, among all possible forms, that is properly human, and this determinate form (ahankara, or `self-consciousness') is inherent in what we term the `mental' faculty, that is, precisely that `internal sense' which is designated in Sanskrit by the name manas, and which is truly the characteristic of human individuality. This faculty ... must be carefully distinguished from pure intellect, which latter, on the contrary, by reason of its universality, must be regarded as existing in all beings and in all states, whatever the modalities through which its existence is manifested,; and one should not see in the `mental' anything beyond what it really is, that is, to use the terminology of logic, a `specific difference' pure and simple, the possession of which does not itself confer on man any effective superiority over other beings.''

Both of these quotes are taken from Guenon's ``The Multiple States of the Being''. If the human state, the individual state and the psychic principle (both mental and sentimental) are all one and the same, how is the individual `asserting himself', `giving meaning and purpose to the world' and `realising himself' not an exaltation of the human state and how are the states Evola describes *beyond* the human state, as opposed to extremities that may be reached within it?

\hfill

\texttt{Cologero on 2011-03-09 at 07:45 said:}

As for ``realising himself'', how is that different from how Guenon describes the ``True Man'', the one who realises all his possibilities? Are you suggesting that the world has no meaning or purpose, that a man acts randomly without meaning, purpose or intent? ``Like, whatever
... '' is not a heroic attitude.


Evola's concern is with the ``I'', which is not individual (``not multipliable''). Compare with Coomaraswamy's essay ``The One and Only Transmigrant''. We've emphasized that there is a difference in perspective between Guenon and Evola. Guenon is strictly a metaphysician, interested in higher states, or really the Transcendent Man, beyond all states. But even the jivanmukti --- the one who has transcended while in this world --- nevertheless continues to act in the world. Is his action random, without purpose, dependent on the external environmental factors of gross and subtle manifestation? Or does he act from the ``center'', the Unmoved Mover? Even the jivanmukti can't avoid ``asserting himself'', just his manner of self-assertion.

Of course, ultimately, there is no ``external environment''. The concept of privation is not a novelty from Evola, even Guenon uses that term.

In Multiple States, Guenon writes about the angelic hierarchy, only to dismiss it. From a metaphysical point of view, all that matters is the existence of multiple states in general, not what they are specifically. However, from a cosmological point of view (and Hermetic, which is a cosmology), these particular states are of great interest.

As for the Individual and the Absolute, Evola writes: ``between the person and universal subject there is a relationship not of coexistence and otherness, rather of continuity and progressivity.'' We will bring that up in the discussion of the One and the Many.

\hfill

\texttt{Aidan on 2011-03-09 at 21:00 said:}

It is frankly astonishing that you can derive any suggestion of an attitude of ``like, whatever'' from my previous comment. Given the reference to individuality being oriented towards and grounded within supra-individual states -- a perspective which transforms both world and individual together -- where on earth do you find even a hint that the world might be without meaning or purpose, or that a man ``acts randomly without meaning, purpose or intent''? When every last part of its order is derived from a still higher order, the world is saturated with meaning and demonstrates perfect purpose; upon the plane of the individual, precisely the same is true of actions which arise directly and spontaneously from the stillness of immovable principle and return to it with complete circularity.

There is no comparison between heroism and divinity. There is not one hero; there are many. There are not many divinities; there is one. Heroism is necessarily attached to individuality while divinity transcends it completely. Compare with Coomaraswamy's `One and Only Transmigrant' indeed! Who is it who transmigrates from incarnation to incarnation? It is not Cologero or Aidan or Evola or Guenon: ``the Lord is the only transmigrant.''

Coomaraswamy and Guenon are at one on this point because both are at one with primordial tradition and their knowledge lacks nothing. Evola, by contrast, is out on his own -- an island in the ocean -- at no point is he talking about the Lord. He means -- as you say -- make YOURSELF God. He means himself, nothing higher: for him, his interior sense is the highest, hence all the sound and fury. It is not a subtle mystery; indeed, it is a crude caricature of true knowledge. Even if he is not oriented towards unity with the Lord, as a brahmin would be (``Whoever is joined unto the Lord is One Spirit'' 1 Corinthians 6:17), he still needs to bow to the Lord, as all kshatriya must, if his deeds are to be genuinely heroic and his heart truly noble.

Yes, the True Man realises himself and all his possibilities but not at the expense of transcendence. He does not have to be transcendent himself -- which is just as well, because he cannot be -- but his relationship to the Transcendent cannot be aberrant if he is to be and remain True.

There would be no problem if there were no insistence that Evola and Guenon are providing related and compatible perspectives. In a fundamental sense, they reveal themselves to be irreconcilable. Either Evola disrupts and deforms Guenon's metaphysics or Guenon severely restricts Evola's ambition. Pressing them together without allowing one to fully dominate the other, as though your own third perspective went further than either, will do a disservice to one or the other and conceal elements in both that need to be seen.

The suggestion that was made earlier that knowledge and action are simply two different spiritual paths is misleading in the extreme. No-one would suggest Evola is a representative of karma marga, a path of selfless action, yet this is the only sense in which action and knowledge can be presented as parallel contrasting paths about which individual choices can be made. Karma marga can be contrasted with jnana marga and bhakti marga without any confusion. The issue of the brahmins and kshatriyas -- and the type of knowledge and type of action they represent -- is not a matter of two different paths that an individual may or may not take. These are social castes, not individual paths. It is the ordering of the world that is at stake. Can there be two world orders simultaneously? This is what Evola disingenuously presents before contradicting himself (and placing action above knowledge) and you echo it. Guenon presents knowledge and action as two aspects of the same order and has knowledge superior to action in the sense that correct and appropriate action proceeds out of knowledge and not vice versa. Both cannot be true; the issue can be fudged, however, which makes it come out in Evola's favour by default. Evola's is a recipe for a kingdom divided against itself and it will not stand -- it will be perpetual all-out war followed by collapse and destruction.

Finally, with regard to the angelic hierarchies, Guenon says that ``the existence of extra-human and supra-human states is of little importance to us'' in the sense that we have ``no motive for occupying ourselves especially with them''. He does not dismiss them at all, speaking of angels as beings in their own right having none of the faculties of an individual order, e.g., reason, but in no way being inferior as a result of this and saying that ``the `angelic' states are properly the supra-individual states of manifestation, that is, those pertaining to non-formal manifestation.''

I intend this to be my last comment on this site. The manner in which my earlier comments have been received has spoken volumes and it would be futile to continue in the same vein. In any event, allow me to thank you for your endeavors. I wish you well.

\hfill

\texttt{Matt on 2011-03-09 at 22:15 said:}

Aidan,

I've felt when reading Evola that he places action above knowledge. He may focus on action a lot more , but I think that isdue to his particular persuasion.

Anyways, hopefully you continue to comment. I think your posts are well-written and contian some good points.

\hfill

\texttt{Matt on 2011-03-09 at 22:16 said:}

Correction: First sentence should be... .I've never*

\hfill

\texttt{Cologero on 2011-03-10 at 08:09 said:}

It is important not to equate ``the courage of one's convictions'' with being possessed by a ``fixed idea''. Although this is a distasteful task, nevertheless I need to clarify some points lest anyone think there is any substance to the objections made to the commentaries on Evola's book. (yes, intelligent objections are welcome.)

The commenter objected to these phrases: `` `asserting himself', `giving meaning and purpose to the world' and `realising himself' ''

Gornahoor intersperses other material with the commentary on Evola for a specific purpose: viz, to situate Evola's ideas within a larger perspective. In that way we can properly understand those phrases, dealing with their actual content rather than nitpicking over ``imperfections of expression.'' Thus:

``Asserting oneself'' needs to be understood as the role of Man as the intermediary between Heaven and Earth; as the Will as intermediary between Providence and Destiny.

If the commenter objects to the perfectly legitimate idea of `giving meaning' how else are we to understand that objection? Is it really preferable that the ``universe'' provides meaning that we passively accept? It is just that attitude that Evola opposes.

``Realising oneself'' is again relatable to Guenon. The individual as Spirit (a non-formal state) actualizes or realizes his potentialities in the World. Again, nothing at all objectionable from a Traditional perspective.

Evola is providing a phenomenological study of the ``I'', that is available to anyone making the effort. He it not writing about the will of Evola, but rather the will of the I. In this he should be understood as taking Shankara's perspective against Ramanuja. It is a legitimate point of view, which Evola develops to its logical conclusion. Again, we pointed out that the concept of privation is common to Aquinas and Guenon. Now, of course, anyone can prefer the position of the ``realist'' to the ``idealist'' (as Evola uses those terms). However, to make the discussion worthwhile, he needs to respond to Evola's arguments against the realist, not to wish that they had never been made.

The idea of the Absolute I, or Brahman, is not unique to Evola; perhaps readers are not so accustomed to dealing with the logical consequences of this view.

It is not so absurd, and subjectivity is not a counter-argument. Go back and read what we have written about the Primordial State, particularly as historically documented by Julian Jaynes. Subjectivity and sense of ego are not necessary or even natural to man. Jaynes documents how at some stage, men were not plagued by such feelings. Instead, they acted together as a unified group. There is a reason for all these posts, they fit together.

Finally, we can mention that there is one Hero, with a thousand faces. Someone once wrote that there are just a few causes of human misery --- illness, poverty, oppression --- and are common to every man. However, the Heroic path is unique to each man. Unlike the path to misery, there is no one way to God, or Enlightenment. Instead each man has to find his own way and deal with his own Destiny. That is why the Hero has a thousand faces.

\hfill

\texttt{Albion on 2011-03-10 at 08:42 said:}

Hi, it's good to be allowed to ask questions, and thank you for this marvelous site.

My question is, if Guenon (I don't have the exact reference) was worried about ``spiritual counterfeits'' and ``amalgams'' that had the form of traditionalism, but whose principle was entirely demonic in its unity, why was it that the ``West'' immediately sprang to mind in the sense that the history of the West for the last thousand years seems to be one entirely of REVOLUTION and re-synthesis which always aspires to conserve? I don't know enough about the East to ask a similar question, but the East seems to me to be lacking in some form of energy. Is this merely my perception as a Westerner under the shadow?

\hfill

\texttt{Graham on 2011-03-10 at 12:04 said:}

I for one cannot understand Aidan's objections. The theory involved is quite easily followed: Evola took the `raw material' of German Idealism and transmuted it into the theosis of a Meister Eckhart or a Plotinus.

\hfill

\texttt{Cologero on 2011-03-10 at 12:09 said:}

I don't understand the point, Albion. You say the west is wracked by REVOLUTION, but where do we see the ``re-synthesis which always aspires to conserve''?

Read the alternatives mentioned by Guenon:
\textit{The Prophecies of Rene Guenon.}\footnote{\url{http://www.gornahoor.net/?p=35}} One of them involves a revolution. Instead of the constant harping about Guenon vs Evola, I wonder if anyone takes any of this seriously? Or is it more fun to re-arrange the deck chairs?

\hfill

\texttt{Graham on 2011-03-10 at 12:30 said:}

Furthermore, if at a *certain stage* in this process one winds up alone like Robinson Crusoe, how can that be construed as an objection to the process itself? Doesn't St. John of the Cross speak of a dark night of the soul? Doesn't Jesus himself ask, on the cross, ``my God, why have you forsaken me?''

\hfill

\texttt{Albion on 2011-03-10 at 20:04 said:}

Thank you for the response. The question was very poorly phrased indeed; I was hoping for some pointers on where Guenon thought the West had gone off the rails (chronologically), and I see you've posted a link -- it just occured to me that all of the revolutions in the West we've glorified (from 1000 or so onwards) are probably suspect as covert ``aspirations'' which are really a betrayal of unity, and therefore, the Good.

As for Evola, I am a neophyte at all this, but clearly if he had thought himself God, then he would not have felt obligated to clearly condemn so many movements, trends, or fads (eg., homosexuality). It's clear that a fair reading of Evola would support a more positive interpretation, which this website appears to be attempting to do.

I support it.

\hfill


\texttt{Albion on 2011-03-14 at 21:30 said:}

``Between these types of existence there is mental activity in the proper sense. In other words: beyond the ideal limit of the reign of pure necessity --- of nature and spontaneity --- and beyond his ``privation'', the individual enjoys in the rational or ideal order of the first level of sufficient actuality and freedom.''

Cologero, could you clarify if the ``first level'' in the last sentence here is referring to the ``second sense'' of will known as power? So thought oscillates between identifying with unreality and privation and the force of power against the self, or fully identifying with its full purpose? Is this what he is concluding the passage by saying?

\hfill

\texttt{Cologero on 2011-03-15 at 00:03 said:}

He is saying that beyond the first stage of privation and spontaneity, the individual first comes to know as ideas, or the ideal. Thus, through thinking, a man comes to see and understand his situation. That is the situation of philosophy in the West, viz, from German Idealism to Giovanni Gentile. But that is insufficient unless and until it becomes a process of self-transformation, can the ideal become real. This is real knowledge, ``carnal knowledge'', when knowing and being are the same.

I recommend ``The World as Power'' by John Woodroffe. This extended work will clarify this rather terse exposition by Evola. (Or else read the 700 pages of the complete system of Magical Idealism.)

\hfill


\texttt{Albion on 2011-03-19 at 01:15 said: }

Since the purpose of the sight is not to defend Evola, but to build on him, then some elective affinity would seem to be in order, just to understand his position. However, Cologero may have stumbled on something here, which goes to the root of the misunderstanding of Platonism as somehow anti-matter. From the higher perspective, ``matter" would of course be entirely good insofar as it existed at all, but would necessarily be ``evil" insofar as it tended to deprive one of that higher perspective. Which is exactly what has happened. Plato never spoke as if we did not have bodies (George Parkin Grant). And one must not overlook the sophistication and symbolism and secret doctrine which operated to imbue statements that were ``wrong" if taken ; context, more than terminology, is everything.


\hfill

\texttt{Albion on 2011-03-19 at 12:59 said: }

I am beginning to think Platonism is nothing like what we have been told that it was. Necessity, without freedom is evil, just as freedom without necessity is evil. The traditionalists (or Platonists) don't get to talk about the divine loop where the end meets the beginning, generally because everyone quarrels with their quest. You can't have ``Western" without Platonism. I don't see where Evola is any more anti-material than Platonism, in fact, far less so.


\hfill

\texttt{Albion on 2011-03-19 at 13:00 said: }

Perhaps one could say that untransfigured matter, or matter rendered untransfigurable, is what is ``evil". Which can be done.


\hfill

\texttt{Will on 2011-03-20 at 11:25 said: }

The five numbered points from Evola's book are strongly parallel to the methodology of Tantric sadhana, so much so that I wonder whether Evola got it from Woodroffe or some other source, or whether he somehow figured it out on his own. This methodology survives to this day in Tantric Buddhist practices in which yidam deities are visualized and meditated upon and prayed to, until they become as real and more real than ordinary reality.

These visualizations are one of the primary uses of thangka paintings: they provide the image to be formed in the mind. There is some evidence that icons in the Christian tradition may have been used in similar ways, and perhaps still are in certain Orthodox circles. In both cases, the painter must go through certain preparatory and purificatory rituals prior to beginning the work. In both cases, the resultant images (if they are successful) are regarded as holy, and as retaining something of the essence of the deity or saint which they portray.


\hfill

\texttt{kadambari on 2011-03-21 at 12:22 said: }

@Cologero

``Idiot" originally referred to a ``person so mentally deficient as to be incapable of ordinary reasoning".

Alain de benoist gives an interesting definition of Idiot. I find his understanding of antiquity of interest, and do not know much about his politics. He writes:

``The notions of citizenship, liberty, or equality of political rights, as well as of popular sovereignty, were intimately interrelated. The most essential element in the notion of citizenship was someone's origin and heritage. Pericles was the ``son of Xanthippus from the deme of Cholargus." Beginning in 451 b.c., one had to be born of an Athenian mother and father in order to become a citizen. Defined by his heritage, the citizen (polítes) is opposed to idiótes, the non-citizen—a designation that quickly took on a pejorative meaning (from the notion of the rootless individual one arrived at the notion of ``idiot"). Citizenship as function derived thus from the notion of citizenship as status, which was the exclusive prerogative of birth. To be a citizen meant, in the fullest sense of the word, to have a homeland, that is, to have both a homeland and a history. One is born an Athenian—one does not become one (with rare exceptions). Furthermore, the Athenian tradition discouraged mixed marriages. Political equality, established by law, flowed from common origins that sanctioned it as well. Only birth conferred individual politeía.

That is essentially how caste was : people with a certain history and way of life…so being out of caste is to have no history and be a rootless individial…

\url{http://home.alphalink.com.au/\~{}radnat/debenoist/alain14.html}


\hfill

\texttt{John Bardis on 2011-03-21 at 16:47 said: }

In \textit{The Yoga of Power} Evola writes:

``According to Indo-Aryan tradition, a caste corresponds to an individual differentiated nature, and to the body of laws (dharma) one should follow. In the Left-Hand Path, however, dharma is seen as a limitation, or as a conditioning to be overcome…Likewise, it does not matter to what caste a person pursuing absolute liberation belongs. Obviously, this is a case of apparent `democracy'. The capability of effectively walking on the path produces differentiations among people that are more clear-cut and real than the ones due to caste membership, especially when the latter have lost their original and legitimate foundation." (page 97)

Perhaps this is what Gurdjieff's toasts to the idiots (or whatever it was called) was all about. I forget how it works, but there were forty or more different types of idiot, and each person had the task of determining which, exactly, kind of idiot he or she was–resulting in the downing of vast quatities of alcohol.


\hfill

\texttt{kadambari on 2011-03-21 at 18:57 said: }

``Likewise, it does not matter to what caste a person pursuing absolute liberation belongs."

Yes in classical India you had the sanyasin (samana) who was beyond caste and religion (you could not demand of a sanyasin what his caste traditionally). However, this did not injure the overall political set up, even the sanyasin is something at the beginning. This is what I find so amazing about this civilization: all these elements could co-exist (the samanas later lead to the rise of Buddhism), which in turn co-esists with what precedes it, so a civilization self-corrects as long as its essential character is not disturbed in the sense that the essential way of life is destroyed; all is self-correction comes to a halt when classical civilization is destroyed by outside forces in the middle ages…So this is a case of pluralism which has none of the aspects we associate with pluralism in the modern sense, it is not multicultural because the differences are a facet of a larger civilization which encompasses it all in a unity….


\hfill

\texttt{John Bardis on 2011-03-21 at 19:37 said: }

So, then, if I may ask–What caste are you?


\hfill

\texttt{kadambari on 2011-03-21 at 20:40 said: }

Well we are Brahmins and retain some memory but today the civilization there is not one that can be properly be called ``Hindu"; so I feel that we are helpless residues of a civilization that harldy exists anymore except among some remnants who carry a memory of it, people who are small in numbers; when I read this blog I feel Westerners are also dealing with similar dilemnas in a Western context…Yes I suppose we are also idiots in Benoist's sense, being unable to relate to or sympathize with the current form of civilization which has taken root there, so rootless in this sense…

I suppose there is not much one can do but our tiny part to keep the memory of these things alive…Imagine if it were not for all these traditionalist writers we would not be concerning ourselves with such questions as a larger group, even if we might do so quietly as individuals! These things have to have impact amongst the public to have a future and cannot just be confined to isolated individuals in academia…A knowledge that there is a problem is the first step…


\hfill

\texttt{John Bardis on 2011-03-22 at 12:18 said: }

In the town where I live, a few miles south of Atlanta, there's a Hindu temple. Or, actually, there are two of them side by side. They are quite magnificent. The main ``god" of the first temple has to be the most amazing statue I've ever seen.

But in another part of Atlanta there is a whole temple complex that makes our two temples seem quite small in comparison.

There was a time in my life when I had the good fortunate to attend a good many concerts of classical Indian music.

Obviously Hinduism (or whatever you call it) isn't going anywhere. It's definitely here to stay.


\hfill

\texttt{kadambari on 2011-03-22 at 14:23 said: }

I hope the Hindus are also doing something for the community… I can say that I myself know many things about our religion from the great Western scholars; the Buddhism book my brother read in college was written by a Catholic priest and it was a very good introduction I recall. Of course there are certain things that are passed down and you know as being a part of the culture, but we are always amazed by the dedicated Westerners who study it to such depths; the German scholars preserved many things when they were dying out even in our regions; genuine quest for knowledge transcends cultures…

I wish Indians would be exposed to the Greeks for they would really love them and learn so much from them and find them not different at all, but what can you do? With the current government run educational system there, Hindus barely understand even their own culture and history…


\hfill

\texttt{kadambari on 2011-03-23 at 08:31 said: }

@John Bardis

The kinds of things you talk about are probably organized by South Indians, who as classical civilization was destroyed much later in the south still have some memory of architecture and other cultrual memories and so on, and can be benign people…The North (I am speaking of places like Delhi) has a thuggish culture because it was totally destroyed at several points and hence there is a barbaric legacy which survives there, I really dislike the culture in those parts, however, the mountian/hill areas in the North where people took refuge during turmoil historically that I am familiar with are completely benign and still very pleasant to live…

As for Indian classical music, most of the great names (about ten) have died, there are only two great masters left…I wonder with interest diminishing in these things when we will again hear music composed by such masters again…


\hfill

\texttt{Yuhomo on 2011-03-23 at 09:54 said: }

Evola is more magician or philosopher than metaphysician. Here he tends to pit freedom and will against necessity, yet he forgets that freedom doesn't exist without necessity. The divine by necessity willed being and cosmos, not the fall. And to be an individual being is not a constraint or weakness since its will is to be an individual whose immanence does not run counter to transcendence. Thus the statement, ``There is no other evil beyond necessity, of which matter, brute existence, is the evidence" is absurd. Then in the moral argument Evola tends to be too Nietzschean in wanting to create an amoral morality based on power. But what do I know?


\hfill

\texttt{John Bardis on 2011-03-23 at 13:25 said: }

Yes, I was told that the temples near where I live were for south Indians.

You will have to forgive the fact that my ignorance on these matters far outweighs my very real interest.

I believe that Buddhism was still prevalent in India, though only in northern India, at the time of the Muslim invasion. And, of course, this invasion only penetrated so far. So I suppose it does make sense that the traditions would be more intact in the south.

As for the music, it seems to me quite similar to western classical music in that it could only be created and can only be performed due to the existence of child prodigies. I believe both kinds of music are still being performed at a very high level. But perhaps creation has come to a stand-still. And this seems to be the case in many fields–literature, philosophy, art. Why this is, I don't know. Perhaps there are times of creation and times of preservation. Or perhaps all the emphasis now is on technology and popular culture.


\hfill

\texttt{kadambari on 2011-03-23 at 15:46 said: }

@John Bardis

Yes the South is quite different from the North, they have their own variety of HInduism and their peculiar languages which sound even peculiar to us and their ways are quite different, but I find in spite of all these differences, the people there are gentler compared to the North; many of them obey rules as compared to the North and are capable of better organization amongst themselves, whereas the North is very unruly and disorderly owing to its bad history. Of course, there are subtleties across regions and your geography will play a part in the peculiar aesthetics of your ethnic group…

Indians unfortunately are not exposed to Western classical music that much, which is a shame because it deprives them of something amazing out there just as they are not exposed to Greco-Roman civilization so cannot properly understand Western civilization and are mostly only exposed to English in terms of language. It is not that there is no interest in these things; they just have not been exposed…There needs to be exposure to these things if Indians can have a proper perspective on civilizations not their own…

As far as Western music is concerned, most are only exposed to pop music; you see the Chinese and Japanese really into Western classical music…I was lucky to have parents who appreciated Western classical music, so was lucky to grow up appreciating both classical Indian and classical Western music…

These days, there is much exposure to books in our region which is a good thing about the internet; today you can find most of the classics on the internet which is good for areas where they are not easily accessible. My husband who is in the sciences recalls when he was in college, the librarian was a Kashmiri Pandit, so in addition to the science books the librarian had a very peculiar section of other books according to his taste such as works by Aurel Stein, the explorer Richard Burton and old British histories of India (which contained history without whitewash) etc., which he says opened him up to understanding Indian history properly, as history is not taught in the schools correctly, as it is all written with a political agenda and dominated by Marxist historians and rewriting of history to serve political ends, the truth is not presented. Most people end up studying technical subjects and the sciences to survive in the modern world and find themselves when they get older woefully ignorant of their own civilization apart from what their parents have taught them and often want a refund on their education…


\hfill

\texttt{kadambari on 2011-03-23 at 17:04 said: }

@John Bardis

We need geniuses to bring back the music, but proper patronage is also required to sustain the arts. The classical musical heritage of the South has to to with the renaissance of the arts during the Vijayanagar empire, the last Hindu stronghold. The rulers patronized the arts and all Brahmanical traditions such as music.

\url{http://www.planetradiocity.com/musicopedia/music\_decade.php?conid=2311}

\url{http://www.vasanthvisuals.com/hampi.html}

There are ragas not which have disappeared due to lack of use, many we cannot even imagine because they have been destroyed with civilization in certain parts, so what is existing cannot continue wihout proper patronage. Unfortunately, the way the Indian State is today controlled by the socialist Congress ever since after independence with the exception of five years, leading to India's current misery and ongoing deracination and overpopulation, I am skeptical of preservation of the arts unless some other party comes into power nationally as the native culture is looked down upon and not encouraged by the state…


\hfill

\texttt{Cologero on 2011-03-23 at 19:25 said: }

The ``divine" will is absolutely free and does not act out of necessity. It the divine is infinite, where would the ``necessity" come from?

Evola is speaking about trans-human states, not the ``individual" man as you seem to mean it.

Evola criticizes Nietzsche's ``morality" because the latter did not understand ``asceticism".

Please read the material carefully prior to commenting.

Evola may be wrong, but not for the reasons you present.

To emphasize once again, we are not out to convert, but rather to elucidate Evola and Guenon, in particular. The only issue at stake is whether or not we have presented Evola's philosophy clearly and accurately.


\hfill

\texttt{Yumo on 2011-03-24 at 09:18 said: }

The divine will is not absolutely free in the sense that you imply. It cannot not will the universe, which occurs as a result of its infinitude. Necessity relates to the absolute, and freedom to the infinite, both of which coincide. If the divine is infinite how could it not also include necessity? But necessity does not relate to impossibility or falsehood as Evola would imply.

I know what Evola is trying to say but his accuracy is off, and for reasons which I suspect to be his early influences, to reduce everything to the freedom of the will. So while those he criticized erred in the direction that negated transcendence, he errs in the direction that negates ethical virtue.

He says that matter is a sign of imperfection as if necessity of willing the universe was not necessary, but then the divine would not be infinite nor free. Matter is a sign of the will to matter which enabled the divine to be immanent and therefore is what it is.


\hfill

\texttt{kadambari on 2011-03-25 at 09:26 said: }

``I believe that Buddhism was still prevalent in India, though only in northern India, at the time of the Muslim invasion. And, of course, this invasion only penetrated so far."

The destruction was massive and total in most areas which never recovered. I recall the European who discovered the ellora caves. The people of the area had become so wretched, he could not imagine that the same people had once produced the art in Ellora….


\hfill

\texttt{Matt on 2011-03-25 at 16:03 said: }

Yumo, 

How does Evola negate ethics in his thought? He stated many times in his writings that he takes virtue into account, but virtue as it was traditionally defined (virtus as ``power", the power to act in the right way in the right situation). Gornahoor has a post on St. Anthony and Evola that can give you a good idea of the ethics in Evola's thought.


\hfill

\texttt{John Bardis on 2011-03-25 at 22:19 said: }

Dear Kadambari,

Thank you for those links about the Vajayanager empire. They give a context to things that I have seen and heard. These things aren't simply from India. They are from a particular place in India at a particular time.

And if these things that I have seen and heard are just the last rays of light from a star that no longer exists–well, is there anything unique in that? Here is something I happened to read today:

``All that is mortal lives in this fear of death; every new birth augments the fear by one new reason, for it augments what is mortal. Without ceasing, the womb of the indefatigable earth gives birth to what is new, each bound to die, each awaiting the day of its journey into darkness with fear and trembling."

And yet we know that this music, this art, this religious devotion, this philosophical insight, will never die.

\hfill

\end{sffamily}
\end{footnotesize}

